\documentclass{article}

\usepackage{fontsize}
\changefontsize{14}

\usepackage[a4paper, top=1.5cm, bottom=1.5cm, left=2.5cm, right=1.5cm]{geometry}

\usepackage[T2A]{fontenc}   % кодировка шрифтов для кириллицы
\usepackage[utf8]{inputenc} % кодировка исходного файла (UTF-8)
\usepackage[english, russian]{babel} % правила переносов и типографики

\usepackage{romannum} % римские цифры

\usepackage{amsmath} % вставки текста в формулы, нумерация формул

\usepackage{siunitx} %значок градуса

\usepackage{makecell} %переносы в таблице



\begin{document}

\pagenumbering{arabic}


%ТИТУЛЬНИК


\thispagestyle{empty}

\begingroup
\setlength{\parindent}{0pt}

\begin{bfseries}

\vspace*{0cm}
Министерство образования и науки РФ

\vspace{2cm}
Московский государственный технический университет имени\\Н.Э. Баумана

\vspace{2cm}
Факультет «Фундаментальные науки»

\vspace{2cm}
Кафедра «Теоретическая механика» имени профессора\\ Н.Е. Жуковского

\vspace{2cm}

Домашнее задание № 2\\
Кинематика плоского движения твердого тела

\end{bfseries}

\vspace{2cm}

\textbf{Вариант} \hspace{0.5cm} 20

\vspace{0.5cm}

\textbf{Группа} \hspace{0.5cm} СМ11-21Б

\vspace{0.5cm}

\textbf{Студент} \hspace{0.5cm} Усачев Е.М. \hspace{1.5cm} \underline{\hspace{3cm}} 
\hspace{1cm}\underline{\hspace{3cm}}\\
\small\hspace*{8.5cm} подпись \hspace{2.7 cm} дата \normalsize

\vspace{0cm}
\textbf{Преподаватель} \hspace{0.5cm} Белый В.С.


\vfill
\textbf{Москва 2026 год}
\vspace{1cm}

\endgroup

%ТИТУЛЬНИК


\newpage

\begin{center}\large \textbf{Кинематика плоского движения твердого тела} \end{center}

\noindent Дано:\\ 
$S_{A_1} = 0,2t^2$ м; $\varphi = t^2$ рад; $r = 0,1$ м; $R = 0,3$ м; $t^* = 1$ с.

\medskip

\noindent Необходимо:

-- для момента времени $t^* = 1$ с определить скорости и ускорения точек $A$, $B$, $C$, угловые скорости и ускорения всех звеньев механизма;

-- найти положение мгновенного центра скоростей (МЦС) звеньев и мгновенного центра ускорений (МЦУ) звена 2, с их помощью проверить правильность нахождения скорости и ускорения точки $B$;

-- нанести на рисунок механизма векторы скоростей и ускорений точек $A$, $B$, $C$, обозначить круговыми стрелками направления угловых скоростей и ускорений звеньев;

\bigskip 
\bigskip 
\bigskip 
\bigskip 
\bigskip
\bigskip
\bigskip
\bigskip
\bigskip
\bigskip
\bigskip
\bigskip
\bigskip
\bigskip
\bigskip 
\bigskip 
\bigskip 
\bigskip 
\bigskip
\bigskip
\bigskip
\bigskip
\bigskip
\bigskip
\bigskip
\bigskip
\bigskip
\bigskip
\bigskip
\bigskip
\bigskip
\bigskip
\begin{center}\text{Рис. 1. Схема механизма} \end{center}


\newpage

\textit{\textsf{Кинематические характеристики движения точек  
и звеньев механизма}}

\bigskip

Сначала определим виды движения звеньев данного механизма: звено 1 (рейка) совершает поступательное движение, звено 2 (шестерня) — плоское движение, звено 3 (рейка) — поступательное, звено 4 (шестерня) — вращательное вокруг оси $O(z)$.

\medskip

Звено 4 (вращательное движение). По заданному закону вращения $\varphi = t^2$ рад определим для момента времени $t^* = 1$ с:

\begin{equation}
\omega_4 = |\dot{\varphi}| = |2t| = 2 \text{ рад/с;}
\end{equation}

\begin{equation}
\varepsilon_4 = |\ddot{\varphi}| = |2| = 2 \text{ рад/с}^2.
\end{equation}
\medskip
Скорость и касательное ускорение точки $E$ (точка контакта звеньев 4 и 3):
\begin{equation}
v_E = \omega_4 r = 2 \cdot 0,1 = 0,2 \text{ м/с;}
\end{equation}
\begin{equation}
a_E^\tau = \varepsilon_4 r = 2 \cdot 0,1 = 0,2 \text{ м/с}^2.
\end{equation}
\medskip
Звено 3 (поступательное движение). Ввиду отсутствия проскальзывания:
\begin{equation}
v_3 = v_E = 0,2 \text{ м/с;} \quad a_3 = a_E^\tau = 0,2 \text{ м/с}^2.
\end{equation}
\medskip
Точка $B$ принадлежит звену 3, поэтому:
\begin{equation}
v_B = v_3 = 0,2 \text{ м/с;} \quad a_B^\tau = a_3 = 0,2 \text{ м/с}^2.
\end{equation}
\medskip
Звено 1 (поступательное движение). По закону $S_{A_1} = 0,2t^2$ м:
\begin{equation}
v_1 = |\dot{S}_{A_1}| = |0,4t| = 0,4 \text{ м/с;}
\end{equation}
\begin{equation}
a_1 = |\ddot{S}_{A_1}| = |0,4| = 0,4 \text{ м/с}^2.
\end{equation}
\medskip
Точка $A$ принадлежит звену 1, поэтому:
\begin{equation}
v_A = v_1 = 0,4 \text{ м/с;} \quad a_A^\tau = a_1 = 0,4 \text{ м/с}^2.
\end{equation}

\medskip

Звено 2 (плоское движение). Рассмотрим движение шестерни 2. Скорости точек $A$ и $B$ известны. Определим угловую скорость звена 2. Точка $C$ — полюс шестерни. 
\begin{equation}
    \vec{v_A} = \vec{v_C} + \vec{v_{AC}} = \vec{v_C} + \vec{\omega_2} \times \vec{R_1} ;
\end{equation}
\begin{equation}
    \vec{v_B} = \vec{v_C} + \vec{v_{BC}} = \vec{v_C} + \vec{\omega_2} \times \vec{R_2}.
\end{equation}
\medskip
Отсюда:
\begin{equation}
    \vec{v_A} - \vec{v_{AC}} = \vec{v_B} - \vec{v_{BC}};
\end{equation}
\begin{equation}
    \vec{v_A} + \vec{v_{BC}} = \vec{v_B} + \vec{v_{AC}}.
\end{equation}
\medskip
Так как вектора $\vec{R_1}$ и $\vec{R_2}$ равны по модулю (см. рис. 2), то равны модули скоростей точек $A$ и $B$ относительно точки $C$:
\begin{equation}
    {v}_{AC}={v}_{BC} = \omega R.
\end{equation}
Спроецируем уравнение $(13)$ на ось $Ox$:
\begin{equation}
    {v_A} - {v_{BC}} = {v_B} + {v_{AC}}.
\end{equation}
\medskip
Тогда с учетом уравнения $(14)$ получаем:
\begin{equation}
v_{AC} = v_{BC} = \frac{v_A - v_B}{2} = \frac{0,4 - 0,2}{2} = 0,1 \text{ м/с}.
\end{equation}
Тогда из уравнения $(14)$:
\begin{equation}
\omega_2 = \frac{v_{BC}}{R} = \frac{0,1}{0,3} = 0,333 \text{ рад/с}.
\end{equation}
Спроецируем выражение $(10)$ на ось $Ox$:
\begin{equation}
    v_A = v_C + v_{AC};
\end{equation}
\begin{equation}
v_C = v_A - v_{AC} = 0,4 - 0,1 = 0,3 \text{ м/с}.
\end{equation}
\medskip
Определим ускорения точек звена 2. Так как точка $C$ - полюс, то:
\begin{equation}
    \vec{a}_A = \vec{a}_C + \vec{a}_{AC} = \vec{a}_C + \vec{a}_{AC}^\tau + \vec{a}_{AC}^n;
\end{equation}
\begin{equation}
    \vec{a}_B = \vec{a}_C + \vec{a}_{BC} = \vec{a}_C + \vec{a}_{BC}^\tau + \vec{a}_{BC}^n;
\end{equation}
\medskip
Известно, что:
\begin{equation}
    \vec{a}_A = \vec{a}_A^\tau + \vec{a}_A^n;
\end{equation}
\begin{equation}
    \vec{a}_B = \vec{a}_B^\tau + \vec{a}_B^n.
\end{equation}
\medskip
Известно, что $\vec{a}_C^n = 0$, так как звено 2 находится между двумя рейками. Из уравнений $(19)$ и $(21)$ получаем:
\begin{equation}
    \vec{a}_A^\tau + \vec{a}_A^n = \vec{a}_C + \vec{a}_{AC}^\tau + \vec{a}_{AC}^n;
\end{equation}
\medskip
Спроецируем уравнение $(23)$ на ось $Oy$:
\begin{equation}
    {a}_A^n = {a}_{AC}^n = \omega_2^2R = 0,33^2 \cdot 0,3 = 0,033 {\text{ м/с}}^2.
\end{equation}
\medskip
Из уравнений $(20)$ и $(22)$ получаем:
\begin{equation}
    \vec{a}_B^\tau + \vec{a}_B^n = \vec{a}_C + \vec{a}_{BC}^\tau + \vec{a}_{BC}^n.
\end{equation}
\medskip
Спроецируем уравнение $(25)$ на ось $Oy$:
\begin{equation}
    -{a}_B^n = -{a}_{BC}^n;
\end{equation}
\begin{equation}
    {a}_B^n = {a}_{BC}^n = \omega_2^2R =0,33^2 \cdot 0,3 =0,033 {\text{ м/с}}^2.
\end{equation}
\medskip
Полные ускорения точек $A$ и $B$:
\begin{equation}
a_A = \sqrt{(a_A^\tau)^2 + (a_A^n)^2} = \sqrt{0,4^2 + 0,033^2} = 0,401 \text{ м/с}^2;
\end{equation}
\begin{equation}
a_B = \sqrt{(a_B^\tau)^2 + (a_B^n)^2} = \sqrt{0,2^2 + 0,033^2} = 0,203 \text{ м/с}^2.
\end{equation}
\medskip
Из уравнений $(23)$ и $(24)$ получаем:
\begin{equation}
    \vec{a}_A^\tau = \vec{a}_C^\tau +\vec{a}_{AC}^\tau.
\end{equation}
\medskip
Спроецируем уравнение $(30)$ на ось $Ox$:
\begin{equation}
    {a}_A^\tau = {a}_C^\tau +{a}_{AC}^\tau.
\end{equation}
\medskip
Из уравнений $(25)$ и $(27)$ получаем:
\begin{equation}
    \vec{a}_B^\tau = \vec{a}_C^\tau +\vec{a}_{BC}^\tau.
\end{equation}
\medskip
Спроецируем уравнение $(30)$ на ось $Ox$:
\begin{equation}
    {a}_B^\tau = {a}_C^\tau - {a}_{BC}^\tau.
\end{equation}
Так как ${a}_{AC}^\tau = {a}_{BC}^\tau =\varepsilon_2R$, то из уравнений $(31)$ и $(33)$ получаем:
\begin{equation}
a_{AC}^\tau = a_{BC}^\tau = \frac{a_A^\tau - a_B^\tau}{2} = \frac{0,4 - 0,2}{2} = 0,1 \text{ м/с}^2.
\end{equation}
\medskip
Угловое ускорение звена 2:
\begin{equation}
\varepsilon_2 = \frac{a_{BC}^\tau}{R} = \frac{0,1}{0,3} = 0,333 \text{ рад/с}^2.
\end{equation}
\medskip
Касательное ускорение точки $C$ из уравнения $(31)$:
\begin{equation}
a_C^\tau = a_A^\tau - a_{AC}^\tau = 0,4 - 0,1 = 0,3 \text{ м/с}^2.
\end{equation}
\medskip
МЦС звена 2 находится из пропорции:
\begin{equation}
    \frac{AP_2}{BP_2} = \frac{v_A}{v_B};
\end{equation}
\begin{equation}
    \frac{AB+BP_2}{BP_2} = \frac{v_A}{v_B};
\end{equation}
\begin{equation}
    \frac{AB}{BP_2} = \frac{v_A}{v_B} - 1 = \frac{v_A-v_B}{v_B};
\end{equation}
\begin{equation}
    BP_2 = \frac{AB \cdot v_B}{v_A-v_B} = \frac{2R \cdot v_B}{v_A-v_B} = \frac{2 \cdot 0,3 \cdot 0,2}{0,4-0,2} = 0,6 \text{ м}.
\end{equation}
\medskip
Проверка скорости точки $B$:
\begin{equation}
v_B = \omega_2 \cdot BP_2 = 0,333\cdot 0,6 = 0,2 \text{ м/с} \quad \text{(верно)}.
\end{equation}
\medskip
МЦУ звена 2 (точка $Q$):
\begin{equation}
AQ = \frac{a_A}{\sqrt{\varepsilon_2^2 + \omega_2^4}} = \frac{\sqrt{0,4^2 + 0,033^2}}{\sqrt{0,333^2 + 0,333^4}} = 1,142 \text{ м};
\end{equation}
\begin{equation}
\tan \alpha = \frac{\varepsilon_2}{\omega_2^2} = \frac{0,333}{0,333^2}=3,003;
\end{equation}
\begin{equation}
    \alpha = \arctan 3,003 = 72^\circ.
\end{equation}
Расстояние $BQ$ и проверка ускорения точки $B$:
\begin{equation}
BQ = \sqrt{AQ^2 + AB^2 - 2 \cdot AQ \cdot AB \cdot \cos(\angle QAB)} = 0,577 \text{ м};
\end{equation}
\begin{equation}
a_B = BQ \sqrt{\varepsilon_2^2 + \omega_2^4} = 0,577 \sqrt{0,333^2 + 0,333^4} = 0,203 \text{ м/с}^2 \quad \text{(верно)}.
\end{equation}

\medskip
Вывод: в результате расчетов были получены значения кинематических характеристик рассматриваемого механизма в момент времени $t^*=1$ с, которые представлены в таблице 1. Для звена 4 МЦС - точка $O$, для МЦС звена 2 - точка $P$, для звеньев 1,3 МЦС не определены, поскольку эти звенья движутся поступательно. Для звена 3 МЦУ - точка $Q$.

\begin{flushright}
  \small\textit{Таблица 1} \hspace*{0.4cm}
\end{flushright}

\vspace*{-0.9cm}

\begin{center}
  Значения кинематических характеристик в момент $t^*=1$с
\end{center}

\begin{tabular}{|c|c|c|c|c|}
  \hline
  \rule{0pt}{1cm} № & \hspace*{0.5cm} Название величины \hspace*{0.5cm} & \makecell{Обозначение\\ величины} & Значение & \makecell{Единицы\\ измерения} \\[0.5cm]
  \hline
  \rule{0pt}{1cm} 1  & \makecell{Угловая скорость\\звена 2} & $\omega_2$ & 0,333 & рад/с \\[0.5cm]
  
  \hline
\end{tabular}
\newpage
\begin{flushright}
  \small\textit{Окончание табл. 1} \hspace*{0.4cm}
\end{flushright}

\begin{tabular}{|c|c|c|c|c|}
     \hline
     \rule{0pt}{1cm} № & \hspace*{0.5cm} Название величины \hspace*{0.5cm} & \makecell{Обозначение\\ величины} & Значение & \makecell{Единицы\\ измерения} \\[0.5cm]
     \hline
     \rule{0pt}{1cm} 2  & \makecell{Угловая скорость\\звена 4} & $\omega_4$ & 2 & рад/с \\[0.5cm]
     \hline
      \rule{0pt}{1cm} 3  & \makecell{Угловое ускорение\\звена 2} & $\varepsilon_2$ & 0,333 & рад/с$^2$ \\[0.5cm]
     \hline
     \rule{0pt}{1cm} 4  & \makecell{Угловое ускорение\\звена 4} & $\varepsilon_4$ & 2 & рад/с$^2$ \\[0.5cm]
      \hline
      \rule{0pt}{1cm} 5  & \makecell{Скорость\\точки $A$} & $v_A$ & 0,4 & м/с \\[0.5cm]
      \hline
      \rule{0pt}{1cm} 6  & \makecell{Скорость\\точки $B$} & $v_B$ & 0,2 & м/с \\[0.5cm]
      \hline
      \rule{0pt}{1cm} 7  & \makecell{Скорость\\точки $C$} & $v_C$ & 0,3 & м/с \\[0.5cm]
      \hline
      \rule{0pt}{1cm} 8  & \makecell{Ускорение\\точки $A$} & $a_A$ & 0,401 & м/с$^2$ \\[0.5cm]
      \hline
      \rule{0pt}{1cm} 9 & \makecell{Ускорение\\точки $B$} & $a_B$ & 0,203 & м/с$^2$ \\[0.5cm]
      \hline
      \rule{0pt}{1cm} 10 & \makecell{Ускорение\\точки $C$} & $a_C^\tau$ & 0,3 & м/с$^2$ \\[0.5cm]
      \hline
\end{tabular}
\end{document}
